% Mélange acido-basique (Liban, 2006)
On dispose d'une solution aqueuse d'acide nitreux $\ce{HNO2_{(aq)}}$
et d'une solution aqueuse de méthanoate de sodium
($\ce{HCOO^{-}_{(aq)} + Na^{+}_{(aq)}}$).

\begin{donnees}
    $\mathrm{pK_{A1}}(\ce{HNO2\ttslash{}NO$_2^-$})=3,3$~; \
    $\mathrm{pK_{A2}}(\ce{HCOOH\ttslash{}HCOO^-})=3,8$.
\end{donnees}

On mélange un même volume $V=\qty{200}{\mL}$ de chacune des deux solutions
précédentes. La quantité de matière d'acide nitreux introduite dans le mélange
est $n_{1}=\qty{4.0e-2}{\mol}$ et celle de méthanoate de sodium
est $n_{2}=\qty{8.0e-2}{\mol}$.

\begin{enumerate}
    \item \label{item:q1}
          Écrire l'équation de la réaction qui se produit lors du mélange entre
          l'acide nitreux et l'ion méthanoate.
    \item Exprimer, puis calculer, le quotient de réaction $Q_{r,i}$ associé à
          cette équation, dans l'état initial du système chimique.
    \item \label{item:q3}
          Exprimer le quotient de réaction dans l'état d'équilibre $Q_{r,\eq}$
          en fonction des constantes d'acidité des couples puis le calculer.
    \item Conclure sur le sens d'évolution de la réaction écrite à la
          question~\ref{item:q1}.
\end{enumerate}

